\section{Quantitative separation across a three-edge cut}
\label{sec:quantitative-splice}

We now give a constructive bound beyond the double-ladder family. Throughout
this section, separating families have the meaning of Definition~\ref{def:hsep}
and contain no repeated cycles. The hypotheses require finite separation
numbers for the factors; they do not assume a solution of Barnette's conjecture.

Let $G_1,G_2$ be vertex-disjoint finite simple cubic graphs. Choose vertices
$v_1,v_2$ with incident edges
\[
 s_{1i}=v_1x_i,\qquad s_{2i}=v_2y_i\qquad (i=1,2,3).
\]
The indexing specifies a bijection between their neighbours. Delete $v_1,v_2$
and add $t_i=x_iy_i$ for $i=1,2,3$. Denote the resulting graph by $G$, and its
three-edge cut by $T=\{t_1,t_2,t_3\}$. Put
$I_j=E(G_j)\setminus\delta(v_j)$ and $n_j=|V(G_j)|$.

Gorsky, Steiner and Wiederrecht prove preservation of the qualitative
$H^{+-}$ property under the corresponding tight-cut decomposition of cubic,
3-connected bipartite graphs \citep[Lemma 4.5(iv) and Theorem 1.7(iv)]{GorskySteinerWiederrecht2023}.
That property is exactly the assertion $\hsep(G)<\infty$. This section
counts explicit glued cycles and then reduces the number of selected pairs.
No priority is asserted for the quantitative bounds, and the qualitative
decomposition theorem is already known. The proof below
is self-contained and needs neither planarity nor bipartiteness.

\subsection{The quantitative gluing lemma}
For separating families $\mathcal F_1,\mathcal F_2$, define their three states
by the omitted incident edge:
\[
 \mathcal A_i=\{C\in\mathcal F_1:s_{1i}\notin C\},\qquad
 \mathcal B_i=\{D\in\mathcal F_2:s_{2i}\notin D\}.
\]
Every Hamiltonian cycle belongs to exactly one state. Write
$a_i=|\mathcal A_i|$ and $b_i=|\mathcal B_i|$.

\begin{lemma}[Quantitative three-edge gluing]\label{lem:splice}
Suppose $\mathcal F_1,\mathcal F_2$ strongly separate the edges of $G_1,G_2$.
For each $i$, choose subfamilies
$\mathcal U_i\subseteq\mathcal A_i$ and
$\mathcal V_i\subseteq\mathcal B_i$ whose edge unions contain $I_1$ and $I_2$,
respectively. Such subfamilies always exist. Set $p_i=|\mathcal U_i|$ and
$q_i=|\mathcal V_i|$. Then
\begin{equation}\label{eq:refined}
 \hsep(G)\ \le\ \sum_{i=1}^3
       \bigl(p_i b_i+a_i q_i-p_iq_i\bigr)
 \ \le\ \sum_{i=1}^3 a_i b_i.
\end{equation}
The first bound is realized by an explicit separating family obtained from
\begin{equation}\label{eq:pairs}
 \bigcup_{i=1}^3
 \left[(\mathcal U_i\times\mathcal B_i)
       \cup(\mathcal A_i\times\mathcal V_i)\right].
\end{equation}
The subfamilies may be chosen so that
\begin{equation}\label{eq:statecovers}
 2\le p_i\le\min\{a_i,n_1/2\},\qquad
 2\le q_i\le\min\{b_i,n_2/2\}.
\end{equation}
In particular, $a_i,b_i\ge2$ for every state.
\end{lemma}

\begin{proof}
\emph{State coverage.}
For every $e\in I_1$ and each $i$, separation of $(e,s_{1i})$ supplies a
cycle in $\mathcal A_i$ containing $e$. Thus $\mathcal A_i$ covers $I_1$;
the same argument applies to $\mathcal B_i$. Each state is nonempty, also
directly by separation of $(s_{1j},s_{1i})$ for $j\ne i$.
A Hamiltonian cycle in $G_1$ contains $n_1-2$ internal edges, whereas
$|I_1|=3n_1/2-3$. It therefore misses exactly $n_1/2-1\ge1$ internal edges.
No single cycle covers $I_1$. Starting with any cycle in $\mathcal A_i$,
add a cycle containing one still-uncovered internal edge until all are covered.
At most $n_1/2-1$ further cycles are needed. No cycle needs to be chosen twice.
This proves the bounds for $p_i$, and the argument for $q_i$ is identical.

\emph{Gluing and projection.}
If $C\in\mathcal A_i$ and $D\in\mathcal B_i$, deleting the splice vertices
from the cycles gives two spanning paths. Both have terminal indices
$\{1,2,3\}\setminus\{i\}$. Join their corresponding ends with the two
cut edges $t_j$, $j\ne i$. The result, denoted $C\star D$, is one Hamiltonian
cycle of $G$. Its restriction to the two sides, together with the omitted
vertices restored, recovers $C,D$. Thus distinct compatible pairs produce
distinct cycles, even across different states.

The family $\mathcal K$ obtained from \eqref{eq:pairs} has exactly the first
sum in \eqref{eq:refined}: in state $i$ the two products intersect in
$\mathcal U_i\times\mathcal V_i$. Since $p_i\le a_i$ and $q_i\le b_i$,
the saving from the full product is
$(a_i-p_i)(b_i-q_i)\ge0$.

\emph{Separation within a side and at the cut.}
Replace each incident edge $s_{ji}$ by $t_i$ and leave every internal edge
unchanged. This embeds the edge set of each factor injectively into $E(G)$.
Every cycle of $\mathcal F_1$ occurs as the first factor of a pair in
\eqref{eq:pairs}, because every $\mathcal V_i$ is nonempty. Similarly, every
cycle of $\mathcal F_2$ occurs as a second factor. Consequently $\mathcal K$
separates any ordered pair lying in the image of one factor. This includes
all pairs involving at least one cut edge.

\emph{Separation across the two sides.}
Let $e\in I_1$ and $f\in I_2$. Some $D\in\mathcal F_2$ avoids $f$:
for example, separate $(s_{21},f)$. Let $D\in\mathcal B_i$.
Because $\mathcal U_i$ covers $I_1$, it contains a cycle $C$ using $e$.
The selected pair $(C,D)\in\mathcal U_i\times\mathcal B_i$ therefore
separates $(e,f)$. For the reverse direction, choose a cycle in
$\mathcal F_1$ avoiding $e$, identify its state $i$, and use
$\mathcal V_i$ to supply a cycle containing $f$. This pair belongs to
$\mathcal A_i\times\mathcal V_i$. All ordered edge pairs are covered.
\end{proof}

\subsection{Consequences and scope}
\begin{corollary}[Projection and a bound using only the two separation numbers]
\label{cor:scalar}
If $h_1=\hsep(G_1)$ and $h_2=\hsep(G_2)$ are finite, then
\begin{equation}\label{eq:scalar}
 \max\{h_1,h_2\}\le\hsep(G)\le(h_1-4)(h_2-4)+8.
\end{equation}
More generally, $\hsep(G)$ is finite if and only if both factor values are finite.
\end{corollary}
\begin{proof}
Any Hamiltonian cycle of $G$ crosses the nonempty proper cut $T$ a positive
even number of times, hence exactly twice. Its restriction to each side is
one spanning path: a further component would be a cycle disconnected from
the Hamiltonian cycle. Closing each path through $v_j$ gives Hamiltonian
cycles in both factors. Projection preserves membership of all corresponding
edges. Projecting a separating family and removing duplicates gives separating
families in both factors, proving the lower bound and the necessity of finiteness.
Lemma~\ref{lem:splice} proves sufficiency.

Choose minimum families in the two factors. Put $x_i=a_i-2\ge0$ and
$y_i=b_i-2\ge0$. Since $\sum_i a_i=h_1$ and $\sum_i b_i=h_2$,
\begin{align*}
 \sum_i a_i b_i
 &=12+2\sum_i x_i+2\sum_i y_i+\sum_i x_i y_i\\
 &\le12+2(h_1-6)+2(h_2-6)+(h_1-6)(h_2-6)\\
 &=(h_1-4)(h_2-4)+8.
\end{align*}
The middle inequality follows from nonnegativity of the omitted cross terms.
\end{proof}

\begin{corollary}[One cube insertion]\label{cor:cube}
If $G^+$ is obtained by splicing a cubic graph $G$ with a cube and
$\hsep(G)<\infty$, then
\[
 \hsep(G)\le\hsep(G^+)\le2\hsep(G).
\]
\end{corollary}
\begin{proof}
Label the cube as $C_4\square K_2$, with bottom cycle $01230$, top cycle
$45674$, and vertical edges $04,15,26,37$. The following six cycles form a
separating family; every edge has a distinct four-element incidence signature.
Each displayed vertex sequence is closed by returning to 0.
\[
\begin{array}{c|l}
\text{omitted edge at 0}&\text{cycles}\\\hline
01&03215674,\quad03762154\\
03&01237654,\quad01562374\\
04&01265473,\quad01547623
\end{array}
\]
There are two cycles in each state. Lemma~\ref{lem:splice}, using a minimum
family in $G$ and the full cube family, gives $2\sum_i a_i=2\hsep(G)$.
Projection gives the lower bound. Cube symmetry allows any splice vertex.
\end{proof}

\paragraph{Application to Barnette graphs.}
If both factors are Barnette graphs, their splice is again a Barnette graph.
Simplicity and cubicity follow from the construction. For bipartiteness, swap
one factor's color names if necessary so that the new joining edges have
oppositely colored endpoints. Planarity follows by placing the vertex-deleted
embeddings in two disks and matching their three boundary terminals. Every
bijection of three terminals is realizable after choosing their cyclic order
and orientation.

For completeness, consider deletion of at most two vertices in the splice.
If one is deleted on each side, each remaining side is connected by
3-connectivity of the factor, and at least one joining edge survives.
If both deleted vertices lie on one side, each component on that side contains
a surviving splice terminal: in the corresponding factor with just those
vertices deleted, every vertex can still reach the splice vertex.
The other side is connected and joins all these components. With fewer
deletions the same argument applies. Thus the splice is 3-connected.

\paragraph{What this does not establish.}
Finiteness of the factor separation numbers is a hypothesis. The lemma does
not supply Hamiltonian cycles for an unknown factor. It gives no reduction
for graphs without a nontrivial three-edge cut, and its scalar bound is not
the conjectured sharp extremal bound $\lfloor n(n+8)/32\rfloor$.
In particular, it does not prove Conjecture~\ref{conj:extremal-continuation}
or Barnette's conjecture.
The smaller state subfamilies need not be minimum edge covers; the construction
uses verified upper bounds and never identifies them with exact $\hsep$ values.


\paragraph{Independent finite checks.}
The constructor and independent standard-library verifier in
\texttt{artifacts/splice\_lemma\_20260917/} check 36 labelled examples, including
all six terminal bijections for five pairs of factors and proper separating
subfamilies for a nested splice. These are not claimed to represent 36
isomorphism classes. The largest example has 30 vertices. For two prisms
$C_8\square K_2$, the natural terminal matching gives a full-product upper
certificate of size 44 and a refined upper certificate of size 40.
Since its order is $2$ modulo $4$, this example lies outside the odd-parameter
double-ladder family. Neither number is asserted to be its exact $\hsep$.
The verifier enumerates Hamiltonian cycles independently via perfect-matching
complements, checks all ordered edge pairs and the two cardinality formulas,
and records certificate hashes. These checks support the construction;
the general proof above does not depend on them.
